\chapter{放射危害评价}

\section{正常工况剂量核算}


\begin{table}[H]
\centering
\caption{机房外公众 / 工作人员年附加有效剂量}
\begin{tabular}{|c|c|c|c|}
\hline
点位 & 年有效剂量 (mSv) & 管控目标 & 评价 \\
\hline
a 准备间 & 8.04×10⁻⁴ & 0.25/5mSv & 远低于限值 \\
\hline
b 狭窄过道 & 2.47×10⁻⁴ & 0.25mSv & 远低于限值 \\
\hline
c 预留加速器区 & 7.91×10⁻⁴ & 5mSv & 远低于限值 \\
\hline
d 控制室 & 1.29×10⁻² & 5mSv & 远低于限值 \\
\hline
e 屋面平台 & 5.47×10⁻⁵ & 0.25mSv & 远低于限值 \\
\hline
f 防护门外 & 9.94×10⁻³ & 0.25mSv & 远低于限值 \\
\hline
\end{tabular}
\end{table}


\textbf{技师摆位附加剂量}：年均约 0.052mSv；叠加控制室剂量后合计\textbf{0.065mSv / 人・年}；叠加原有岗位历史剂量最高 0.2mSv，合计远低于 5mSv 内控目标、20mSv 国标限值。
\textbf{周边公众最大年剂量 0.00994mSv}，远低于 0.25mSv 内控。

\section{事故工况风险及防控}

\begin{enumerate}
\item \textbf{卡源故障}：定期校验门机联锁、急停、回源系统；
\item \textbf{源丢失被盗}：废源原厂闭环回收，医院不长期贮存废源；
\item \textbf{人员误入}：门禁 + 指示灯 + 视频监控，治疗前清点人员；
\item \textbf{计划失误}：双核对治疗处方、加强人员培训考核；
\item \textbf{屏蔽老化泄漏}：每月墙外剂量巡检。
\end{enumerate}

\section{小结}

正常工况人员、公众受照剂量远低于管控目标；各类事故均有针对性防控措施，潜在照射可控。
